% Generated draft: requires mathematical and source review.
\documentclass[11pt]{article}
\usepackage[margin=25mm]{geometry}
\usepackage{fontspec}
\setmainfont{Latin Modern Roman}
\usepackage{amsmath,amssymb,mathtools,unicode-math}
\setmathfont{Latin Modern Math}
\usepackage{enumitem,needspace,xurl,hyperref}
\setlength{\parindent}{0pt}
\setlength{\parskip}{5pt}
\setlength{\emergencystretch}{3em}
\newcommand{\N}{\mathbb{N}}
\newcommand{\Z}{\mathbb{Z}}
\newcommand{\Q}{\mathbb{Q}}
\newcommand{\R}{\mathbb{R}}
\newcommand{\C}{\mathbb{C}}
\newcommand{\F}{\mathbb{F}}
\newcommand{\A}{\mathbb{A}}
\newcommand{\PP}{\mathbb{P}}
\newcommand{\E}{\mathbb{E}}
\newcommand{\eps}{\varepsilon}
\newcommand{\dd}{\,\mathrm{d}}
\begin{document}
\section*{UnsolvedMath ID: 30002179}
\textbf{Source problem code:} catalogue entry
\subsection{Definitions and mathematical statement}
For a group \(H\) and a field \(K\), define its first Betti number by
\[
b_1(H;K)=\dim_K H_1(H;K),
\]
where group homology uses the trivial \(H\)-action on \(K\). For a finite-index subgroup \(H\leq G\), write \([G:H]\) for its index.

The source uses the phrase ``inverse system of finite-index normal subgroups'' without specifying its indexing order or transition maps. Under the usual descending-family interpretation, let \(I\) be a nonempty directed partially ordered set and let \((G_i)_{i\in I}\) satisfy
\[
G_i\trianglelefteq G,\qquad [G:G_i]<\infty,\qquad
 i\leq j\Longrightarrow G_j\subseteq G_i,\qquad
\bigcap_{i\in I}G_i=\{1\}.
\]
The transition maps are then the subgroup inclusions, and limits over \(I\) are limits of nets. The existence of such a family implies that \(G\) is residually finite: every nonidentity element is excluded by some finite-index normal subgroup.

\textbf{Question, under this interpretation.} For every finitely presented residually finite group \(G\), does there exist a number \(L(G)\in\R_{\geq 0}\) such that, for every field \(K\) and every family \((G_i)_{i\in I}\) as above,
\[
\lim_{i\in I}\frac{b_1(G_i;K)}{[G:G_i]}=L(G)?
\]
Explicitly, this asks that for every \(\varepsilon>0\) there exist \(i_0\in I\) such that
\[
i\geq i_0\Longrightarrow
\left|\frac{b_1(G_i;K)}{[G:G_i]}-L(G)\right|<\varepsilon.
\]
Thus the asserted value may depend on \(G\), but must be independent of both the subgroup system and the field, including its characteristic. The original wording quantifies over all finitely presented groups; for groups admitting no such system, its universal quantification over systems is vacuous.

The directed-family convention above requires source confirmation; the record does not settle whether only descending sequences were intended. The supplied literature assessment discusses convergence and subgroup-chain independence in positive characteristic, while the supplied statement additionally requires independence of the field. The full statement is retained here. Its supplied open-status label and review labels are curation claims, not independently verified mathematical conclusions.
\subsection{Short English statement}
For a finitely presented group admitting a descending system of finite-index normal subgroups with trivial intersection, ask whether the first Betti numbers divided by subgroup index converge to a value depending only on the group, for every field and every such system. The source leaves the precise meaning of ``inverse system'' unspecified. Its full question includes independence of the field, although its supplied literature summary emphasizes subgroup-chain independence.
\subsection{Sources}
\begin{itemize}
\item Local registry: \texttt{data/unsolvedmath/problems.json}, with \texttt{catalog/problems.json} as the supplementary registry; permanent UnsolvedMath ID \(30002179\), exact source problem code \texttt{catalogue entry}, title ``Chain-Independent Homology Growth in Finite Covers.'' This draft uses the supplied complete record; no registry files were independently inspected.
\item Supplied report attribution: \emph{Low-Dimensional Topology and Number Theory (2013)}, Oberwolfach Reports, pp.~2541--2596, DOI \texttt{10.4171/owr/2012/42}. \url{https://doi.org/10.4171/owr/2012/42}. The record gives proposed year \(2012\) and names no individual proposer.
\item Supplied literature source: ``Oberwolfach Report 42, Betti-limit conjecture.'' \url{https://doi.org/10.4171/owr/2012/42}.
\item Supplied literature source: ``Linnell--Lueck--Sauer, Fp-Betti limits.'' \url{https://arxiv.org/abs/1003.0434}. This is the supplied reference label; no specific theorem or passage is identified in the record.
\item Literature assessment checked at \texttt{2026-08-21}: the supplied curation labels the problem open and states that positive-characteristic convergence and chain independence remain open in general. This is supplied dated curation, not an independently verified literature assessment. The background expressly requests independent verification of time-sensitive status and cited claims.
\item The record labels the statement \texttt{corrected\_verified} and reports manual semantic, typographic, and LaTeX review of its cleanup. These are supplied review claims; this statement-only draft still requires human mathematical review.
\end{itemize}
\end{document}
